% =====================================================================
\chapter{Considerações finais}\label{cap:consideracoes}
% =====================================================================

Este trabalho teve como objetivo comparar o nível e a composição do endividamento
de empresas familiares e não familiares do setor de saúde da B3 entre 2018 e 2025.
A partir da perspectiva da riqueza socioemocional, formulou-se a hipótese de que as
empresas de controle familiar apresentariam menor alavancagem financeira, em razão
da aversão da família à interferência de credores e ao risco de perda de controle.
Para testá-la, 16 companhias foram classificadas segundo o critério de
\citeonline{santos2018}, e as diferenças entre os grupos foram avaliadas por
indicadores contábeis de endividamento, com os testes t de Welch e de
Mann-Whitney, para o período completo e para subperíodos definidos em torno da
pandemia de COVID-19.

Os resultados indicam que as empresas familiares apresentaram, na média de
2018--2025, menor endividamento oneroso que as não familiares, com distância maior
entre 2021 e 2023 e grupos próximos no início e no fim do período. Contudo, as
diferenças não se mostraram estatisticamente significativas na amostra completa,
em nenhum dos indicadores, testes ou subperíodos, de modo que a hipótese H1 não
pode ser confirmada. O endividamento total, por sua vez, foi semelhante nos dois
grupos: descritivamente, as familiares tiveram menos dívida financeira, mas passivo
total equivalente. A composição da dívida por prazo, examinada de forma
exploratória, não apresentou diferença consistente entre os grupos.

As análises de robustez reforçam a cautela na leitura desses resultados. Na amostra
completa, a diferença no endividamento oneroso só foi significativa com a exclusão
da única empresa sem dívida financeira, inverteu o sinal quando se exigiu uma
participação familiar maior para classificar a empresa como familiar e praticamente
desapareceu sem as empresas de comércio e distribuição. Quando a comparação se
restringiu às empresas não reguladas pela ANS e foi feita ano a ano, as familiares
tiveram menos dívida onerosa em todos os exercícios, com diferenças significativas a
5\% em 2022 e 2023; na amostra sem operadoras e sem comércio, a diferença foi
significativa a 5\% apenas em 2022. Esses resultados, porém, dependem de um grupo
não familiar de três empresas, aparecem quase sempre em um único teste e não se
mantêm quando também se excluem os exercícios com prejuízo.

Em resposta ao problema de pesquisa, conclui-se que não foram encontradas
evidências estatisticamente significativas e robustas de que a estrutura de
controle familiar esteja associada a um nível de endividamento diferente nas
companhias abertas do setor de saúde. A evidência descritiva, de caráter
exploratório, sugere que, entre empresas de atividades comparáveis, as familiares
tenderam a usar menos dívida onerosa, padrão compatível com a aversão da família à
interferência de credores, mas também com diferenças de modelo de negócio entre os
subsetores. A ausência de evidência estatística, com estimativas imprecisas, é
coerente com a literatura inconclusiva sobre o tema e com o fato de que a
preservação do controle pode tanto desestimular a dívida quanto levar a família a
preferi-la à emissão de novas ações \cite{gonzalez2013,poletti2022,kayo2018}.

Como contribuição, o trabalho reúne evidências recentes, que incluem o período da
pandemia, sobre um setor que nenhum dos estudos revisados analisa isoladamente, e
mostra que os resultados dependem do critério adotado para
definir empresa familiar, aspecto em que a literatura nacional diverge. Do ponto de
vista prático, os resultados indicam a credores e investidores que a análise do
endividamento de empresas familiares deve considerar o subsetor de atuação e o
regime regulatório das empresas comparadas, e que a classificação de uma companhia como
familiar depende do critério adotado, o que deve ser observado ao comparar estudos.

O estudo apresenta limitações que devem ser consideradas na interpretação dos
resultados. A amostra é pequena, com 16 empresas de um único setor, o que reduz o
poder dos testes estatísticos, torna as estimativas imprecisas e impede
generalizar as conclusões para outros setores. A exigência de dados em todo o
período excluiu as empresas sem informações na base desde 2018, em geral as que
abriram capital mais tarde, o que pode introduzir viés de seleção. Uma das empresas não familiares, a Bradsaúde, não
possui dívida financeira, e sua presença altera o resultado do indicador
principal. A classificação das empresas foi feita com base no Formulário de
Referência mais recente e mantida fixa em todo o período, embora algumas
companhias tenham passado por mudanças relevantes de controle; eventuais erros de
classificação tendem a aproximar os grupos e a atenuar as diferenças. Não há
empresas não familiares no subsetor de comércio e distribuição, e as maiores
empresas da amostra são familiares, de modo que parte das diferenças pode refletir
o subsetor e o tamanho, e não apenas a estrutura de controle. As
operadoras de planos de saúde seguem regras próprias de capital e de provisões
técnicas fixadas pela ANS, e a sua retirada, junto com a do comércio, deixa
apenas três empresas não familiares na amostra mais homogênea.
Aberturas de capital e fusões ocorridas no período alteraram o balanço de várias
empresas e não foram tratadas separadamente, assim como não foi verificado, nas
notas explicativas, se variações pontuais da dívida de curto prazo decorrem de
vencimentos concentrados ou de reclassificações. Também não foi examinado o efeito
da adoção da norma de arrendamentos (CPC 06 (R2)/IFRS 16), em 2019, sobre os
indicadores; a dívida bruta da maioria das empresas com muitos imóveis alugados não
apresenta salto no início de 2019, o que sugere que a Economática não inclui os
passivos de arrendamento nesse indicador, mas a definição não foi confirmada na
documentação da base.
Por fim, o subperíodo anterior à pandemia abrange apenas dois exercícios, e a
análise por subperíodos e a do prazo da dívida, por terem sido definidas durante a
pesquisa, têm caráter exploratório.

Para pesquisas futuras, sugere-se ampliar a amostra para outros setores ou para o
conjunto das companhias não financeiras da B3, o que permitiria estimar modelos de
regressão em painel com controles de tamanho, rentabilidade, tangibilidade e setor.
Sugere-se também testar, como hipótese formulada de antemão, a relação entre o
controle familiar e o prazo da dívida, que neste trabalho foi examinada apenas de
forma exploratória e sem resultado consistente, e examinar se a relação entre a participação da família e o
endividamento é não linear, como apontado por \citeonline{poletti2022}. Outra
possibilidade é classificar as empresas ano a ano, captando as mudanças de controle
ao longo do tempo.
